\chapter{urllib}
\section{简介}
\codeinline{urllib}是Python的一个标准库，它提供了一系列用于操作URL的功能。它包含以下四个模块：
\begin{enumerate}[label=\arabic*)]
    \item \codeinline{urllib.request}：用于打开和读取URL。
    \item \codeinline{urllib.error}：包含了\codeinline{urllib.request}模块中产生的异常。
    \item \codeinline{urllib.parse}：用于解析URL。
    \item \codeinline{urllib.robotparser}：用于解析robots.txt文件。
\end{enumerate}

\section{urllib.request}
\codeinline{urllib.request}模块提供了打开和读取URL的方法 

\begin{lstlisting}[style=demostyle,numbers=none]
# 错误导入
import urllib 
urllib.request

# 正确导入
from urllib import request

urllib.request.urlopen(url, data=None, [timeout, ]*, cafile=None, capath=None, cadefault=False, context=None)
'''
url: URL地址
data: POST请求的数据
timeout: 超时时间
cafile: CA证书
capath: CA证书路径
cadefault: 是否使用默认的CA证书
context: SSL上下文
'''
\end{lstlisting}

\begin{lstlisting}[style=demostyle,numbers=none]
from urllib import request

url = 'http://www.baidu.com'   
response = request.urlopen(url,timeout=1) # 打开URL,timeout=1表示超时时间为1s
print(response.read().decode('utf-8'))    # 读取URL内容并解码

# 获取状态码
print(response.status)
\end{lstlisting}

\subsection{Get}

\href{https://httpbin.org/}{httpbin}是一个HTTP测试网站，可以用来测试HTTP请求，如下是一个简单的例子：
\codeinline{Get}请求:是从指定的资源请求数据。
\begin{lstlisting}[style=demostyle,numbers=none]
from urllib import request

url = 'https://httpbin.org/get'
response = request.urlopen(url,timeout=1) # 打开URL,timeout=1表示超时时间为1s
print(response.read().decode('utf-8'))    # 读取URL内容并解码
\end{lstlisting}

\subsection{Post}
\codeinline{Post}请求:向指定的资源提交数据进行处理请求（例如提交表单或上传文件）。数据被包含在请求体中。POST请求可能会导致新的资源的建立和/或已有资源的修改。

\begin{lstlisting}[style=demostyle,numbers=none]
from urllib import request, parse

url = 'https://httpbin.org/post'
data = bytes(parse.urlencode({'word': 'hello'}), encoding='utf8')
response = request.urlopen(url, data=data, timeout=1)
print(response.read().decode('utf-8'))
\end{lstlisting}

\subsection{Http Headers}
\codeinline{Http Headers}是HTTP请求和响应的头部信息，可以通过\codeinline{add\_header}方法添加头部信息。

\begin{lstlisting}[style=demostyle,numbers=none]
from urllib import request

url = 'https://httpbin.org/get'
headers = {
    'User-Agent': 'Mozilla/4.0 (compatible; MSIE 5.5; Windows NT)',
    'Host': 'httpbin.org',
    'Accept': '*/*',
    'Accept-Language': 'zh-CN,zh;q=0.9',
    'Accept-Encoding': 'gzip, deflate',
    'Connection': 'keep-alive',
    'Cookie': '_gauth=1; _ga=GA1.2.123456789.1599999999; _gid=GA1.2.123456789.1599999999'
}

# Create the Request object with the custom headers
req = request.Request(url, headers=headers)

# Send the request and get the response
response = request.urlopen(req)

# Read and print the response
print(response.read().decode('utf-8'))
\end{lstlisting}

\subsection{Proxy}
\codeinline{Proxy}是一种网络服务，它充当客户端和其他服务器之间的中间服务器。客户端发送请求到代理服务器，代理服务器再发送请求到其他服务器，并将其他服务器的响应返回给客户端。

\begin{lstlisting}[style=demostyle,numbers=none]
import requests

url = 'https://httpbin.org/get'
headers = {
    'User-Agent': 'Mozilla/4.0 (compatible; MSIE 5.5; Windows NT)',
    'Host': 'httpbin.org',
    'Accept': '*/*',
    'Accept-Language': 'zh-CN,zh;q=0.9',
    'Accept-Encoding': 'gzip, deflate',
    'Connection': 'keep-alive',
    'Cookie': '_gauth=1; _ga=GA1.2.123456789.1599999999; _gid=GA1.2.123456789.1599999999'
}

# Define your proxy
proxies = {
    'http': 'http://your_proxy_ip:your_proxy_port',
    'https': 'http://your_proxy_ip:your_proxy_port'
}

# Send the GET request with the headers and proxy
response = requests.get(url, headers=headers, proxies=proxies)

# Print the response
print(response.text)
\end{lstlisting}

\section{requests}
\codeinline{requests}是一个Python的第三方库，它提供了更多的功能和更简单的API，相比于\codeinline{urllib}。
\subsection{安装}
\begin{lstlisting}[style=demostyle,numbers=none]
# 安装
conda install requests
pip install requests
\end{lstlisting}

\subsection{Get}
\begin{lstlisting}[style=demostyle,numbers=none]
import requests

url = 'https://httpbin.org/get'
data = {
    'word': 'hello'
}
response = requests.get(url, params=data)
print(response.text)
\end{lstlisting}

\subsection{Post}
\begin{lstlisting}[style=demostyle,numbers=none]
import requests

url = 'https://httpbin.org/post'
data = {
    'word': 'hello'
}
response = requests.post(url, data=data)
print(response.text)
\end{lstlisting}


